\chapter{El puente entre Rust y QML (cxx-qt)}
\label{cap:puente}

QML sabe hablar con objetos de Qt (\cod{QObject}), que son C++. Rust no es C++. cxx-qt
resuelve el problema: a partir de una descripción en Rust genera el código C++ necesario para
que un \cod{struct} de Rust \emph{sea} un \cod{QObject}, con propiedades, señales y funciones
que QML puede usar.

\section{Anatomía de un objeto puente}

Así es, recortado, el de \cod{Account} (\archivo{src/account.rs}). Los cinco objetos puente
del proyecto siguen exactamente esta forma.

\begin{rust}
#[cxx_qt::bridge]                                   // (1) esto es un puente
pub mod qobject {
    unsafe extern "C++" {                           // (2) tipos de C++ que vamos a usar
        include!("cxx-qt-lib/qstring.h");
        type QString = cxx_qt_lib::QString;
    }

    #[auto_cxx_name]                                // (3) snake_case → camelCase en QML
    extern "RustQt" {
        #[qobject]                                  // (4) un QObject...
        #[qml_element]                              //     ...registrado en QML...
        #[qml_singleton]                            //     ...como singleton
        #[qproperty(bool, signed_in)]               // (5) propiedades
        #[qproperty(bool, busy)]
        #[qproperty(QString, name)]
        #[qproperty(QString, message)]
        #[qproperty(bool, web_login)]
        type Account = super::AccountRust;          // (6) los datos viven en AccountRust

        #[qinvokable]                               // (7) funciones llamables desde QML
        fn sign_in_web(self: Pin<&mut Account>);
        #[qinvokable]
        fn sign_out(self: Pin<&mut Account>);
    }

    impl cxx_qt::Threading for Account {}           // (8) permite qt_thread()
}

pub struct AccountRust {                            // (9) un campo por propiedad, mismo nombre
    signed_in: bool,
    busy: bool,
    name: QString,
    message: QString,
    web_login: bool,
}

impl Default for AccountRust {                      // (10) estado inicial
    fn default() -> Self { /* lee la sesión guardada... */ }
}

impl qobject::Account {                             // (11) el cuerpo de las funciones
    fn sign_out(mut self: Pin<&mut Self>) {
        forget();
        self.as_mut().set_name(QString::default());
        self.as_mut().set_signed_in(false);
    }
}
\end{rust}

\begin{enumerate}
  \item \cod{#[cxx_qt::bridge]} marca el módulo que cxx-qt procesa al compilar.
  \item Dentro de \cod{extern "C++"} se declaran los tipos de C++ que se usarán.
  \cod{QString} es el texto de Qt; para pasar de \cod{String} a \cod{QString} se usa
  \cod{QString::from(&texto)} y al revés \cod{qstring.to_string()}.
  \item \cod{#[auto_cxx_name]}: los nombres de Rust (\cod{sign_in_web}, \cod{signed_in}) se
  ven en QML en camelCase (\cod{signInWeb}, \cod{signedIn}).
  \item \cod{#[qobject]}, \cod{#[qml_element]}, \cod{#[qml_singleton]}: un objeto de Qt,
  visible en QML, del que hay uno solo y se usa por su nombre (\cod{Account}).
  \item \cod{#[qproperty(tipo, nombre)]}: una propiedad. cxx-qt genera para cada una:
  \begin{itemize}
    \item \cod{nombre()} para leerla desde Rust (devuelve una referencia: \cod{*self.busy()}),
    \item \cod{set_nombre(valor)} para escribirla, que además \textbf{avisa a QML} con la señal
    \cod{nombreChanged},
    \item y la propiedad \cod{Account.nombre} en QML, con su notificación.
  \end{itemize}
  \item \cod{type Account = super::AccountRust}: el objeto de Qt guarda dentro un
  \cod{AccountRust}.
  \item \cod{#[qinvokable]}: funciones que QML puede llamar. Aquí solo se \emph{declaran}.
  \item \cod{impl cxx_qt::Threading}: habilita \cod{qt_thread()}, para volver al hilo de Qt
  desde otro hilo.
  \item El struct de datos: \textbf{un campo por cada propiedad, con el mismo nombre y tipo}.
  Puede tener más campos (privados, que QML no ve), como los contadores de peticiones.
  \item \cod{Default} da el estado inicial al crearse el objeto (al arrancar la app).
  \item El cuerpo de las funciones, en un bloque \cod{impl qobject::Account} \textbf{fuera}
  del puente.
\end{enumerate}

\begin{trampa}[«Unsupported verbatim input in ForeignItem»]
Si escribes el cuerpo de una función dentro del bloque \cod{extern "RustQt"}, el generador
falla con este mensaje. Dentro del puente solo van las \emph{firmas}; los cuerpos, en
\cod{impl qobject::X}.
\end{trampa}

\section{\texttt{Pin<\&mut Self>}, \texttt{as\_mut()} y \texttt{rust\_mut()}}

Los objetos de Qt no se pueden mover en memoria (Qt guarda punteros a ellos). Rust lo
expresa con \cod{Pin}: un \cod{Pin<&mut Account>} es «una referencia modificable a un Account
que no se moverá». Tres cosas prácticas:

\begin{description}[leftmargin=1.2em]
  \item[\cod{self.as_mut()}] Para llamar a varios métodos que consumen \cod{Pin<&mut Self>}
  seguidos. Sin él, la primera llamada «se quedaría» con \cod{self}:
\begin{rust}
self.as_mut().set_busy(true);
self.as_mut().set_message(QString::from("Conectando…"));
self.set_name(QString::default());            // la última puede usar self directamente
\end{rust}
  \item[\cod{self.campo}] Leer un campo privado del struct (los que no son propiedades) funciona
  directamente: \cod{self.tracks}, \cod{self.request}.
  \item[\cod{self.as_mut().rust_mut()}] Para \emph{modificar} un campo privado:
\begin{rust}
self.as_mut().rust_mut().request = request;
\end{rust}
\end{description}

En la práctica: copia la forma de las funciones vecinas del mismo archivo. Es mecánico.

\section{Volver al hilo de Qt}

Las propiedades solo se pueden tocar desde el hilo de Qt. El patrón completo, de
\cod{Library::search}:

\begin{rust}
fn search(mut self: Pin<&mut Self>, query: &QString) {
    let query = query.to_string();                     // QString → String (para el hilo)
    let request = self.search_request + 1;
    self.as_mut().rust_mut().search_request = request;
    self.as_mut().set_searching(true);

    let qt = self.qt_thread();                         // (1) el «buzón» del hilo de Qt
    thread::spawn(move || {                            // (2) trabajo lento en otro hilo
        let result = ytm::search(&query);
        let _ = qt.queue(move |mut library| {          // (3) esto corre en el hilo de Qt
            if library.search_request != request {     // (4) ¿sigue siendo la última?
                return;
            }
            library.as_mut().set_searching(false);
            match result {
                Ok(zones) => library.as_mut().set_busqueda(QString::from(&to_json(&zones))),
                Err(error) => library.as_mut().set_error(QString::from(&error)),
            }
        });
    });
}
\end{rust}

\begin{enumerate}
  \item \cod{qt_thread()} devuelve un objeto que se puede mandar a otros hilos.
  \item El hilo hace el trabajo lento. No puede tocar \cod{self}: no lo tiene (y el compilador
  no le dejaría).
  \item \cod{qt.queue(clausura)} pide a Qt que ejecute la clausura en su hilo, pasándole el
  objeto (\cod{library}). Ahí ya se pueden cambiar propiedades.
  \item El contador descarta respuestas que llegan tarde.
\end{enumerate}

\section{Del lado de QML}

\begin{qml}
import app.musicplayer

Text {
    text: Account.signedIn ? Account.name : "Sin sesión"   // lee propiedades
}
PillButton {
    text: "Cerrar sesión"
    onClicked: Account.signOut()                           // llama a una función
}
Connections {                                              // reacciona a un cambio
    target: Account
    function onSignedInChanged() { Library.load() }
}
\end{qml}

\section{Las demás vías entre C++ y Rust}

No todo pasa por cxx-qt. Tres funciones C++ propias y una de Objective-C++ se llaman como
funciones C normales (\cod{extern "C"}), porque no son objetos de Qt sino utilidades:

\begin{center}
\begin{tabularx}{\linewidth}{@{}lX@{}}
\toprule
Función & Qué hace \\
\midrule
\cod{musicplayer_configure_network()} & Desactiva el proxy automático del sistema en Qt (\archivo{src/window.cpp}) \\
\cod{musicplayer_tune_windows()} & Libera la memoria gráfica de las ventanas ocultas o minimizadas; programa los clics de depuración \\
\cod{musicplayer_login(ruta)} & Ventana de inicio de sesión de macOS (\archivo{src/login\_mac.mm}) \\
\cod{musicplayer_login_linux(ruta)} & Ventana de Linux (\archivo{login/src/login\_linux.cpp}) \\
\bottomrule
\end{tabularx}
\end{center}

\section{Resumen: añadir algo al puente}

\begin{enumerate}
  \item Propiedad: \cod{#[qproperty(tipo, nombre)]} en el puente + campo en el struct + valor
  inicial en \cod{Default}.
  \item Función: firma con \cod{#[qinvokable]} en el puente + cuerpo en \cod{impl qobject::X}.
  \item Objeto nuevo: archivo nuevo con su puente, añadirlo a \cod{.files([...])} de
  \archivo{build.rs} y \cod{mod} en \archivo{src/main.rs}.
  \item En QML: \cod{X.nombreEnCamelCase}.
\end{enumerate}
